\documentclass[12pt, spanish]{article}
\usepackage[utf8]{inputenc}
\usepackage{geometry}
\usepackage{graphicx}
\usepackage{enumitem}
\usepackage{multirow}
\usepackage{booktabs}
\usepackage[spanish]{babel}
\usepackage{hyperref}

\geometry{a4paper, margin=1in}

\title{Proyecto - Visualización de Datos}
\date{}

\begin{document}

\maketitle

\section*{Descripción}
El objetivo del proyecto es aplicar los conceptos y herramientas de visualización de datos para resolver una necesidad real de análisis en una empresa u organización. Cada grupo deberá seleccionar una entidad de cualquier sector (salud, educación, economía, seguridad, gobierno, finanzas, etc.) que requiera desarrollar un proyecto de visualización para apoyar el análisis y la toma de decisiones.

\section*{Objetivos del proyecto}
\begin{itemize}
    \item Aplicar los conceptos aprendidos de visualización de datos en un caso práctico de empresa u organización.
    \item Analizar un contexto real para identificar problemas y oportunidades donde la analítica y la visualización aporten valor a la toma de decisiones.
    \item Formular objetivos de visualización y analítica claros y alineados con los objetivos de la empresa u organización.
    \item Integrar y preparar datasets de diferentes fuentes para su uso en proyectos de visualización.
    \item Diseñar y organizar la información de forma clara, estructurada y efectiva para comunicar hallazgos.
    \item Construir dashboards dinámicos y reportes ejecutivos que respondan a preguntas clave.
    \item Comunicar resultados mediante storytelling, desarrollando una narrativa convincente que guíe a la audiencia desde el problema hasta las conclusiones y recomendaciones.
\end{itemize}

\section*{Tareas}
\begin{enumerate}
    \item Comprender el negocio y su contexto: analizar la organización seleccionada y el sector al que pertenece.
    \item Definir el problema y las preguntas generadoras del proyecto, en relación con los objetivos estratégicos de la organización.
    \item Explorar fuentes de datos (Kaggle, Datos Abiertos Colombia, Medellín, Bogotá, entre otras) y recolectar los datasets necesarios.
    \item Preparar y limpiar los datos: depuración, transformación y verificación de calidad.
    \item Modelar los datos: elaborar el modelo lógico de la información (ej. diagrama entidad–relación, modelo dimensional).
    \item Plantear los objetivos de analítica y visualización, definiendo indicadores que respondan a las preguntas del proyecto.
    \item Diseñar la organización de la información y los elementos visuales, aplicando principios de diseño, organización y comunicación clara.
    \item Construir la narrativa (storytelling), que guíe a la audiencia desde el problema hasta las conclusiones y recomendaciones.
    \item Diseñar e implementar los reportes y dashboards dinámicos que integren indicadores, diseño y visualizaciones interactivas.
    \item Presentar los hallazgos finales mediante reportes ejecutivos y dashboards, contando una historia clara y orientada a la toma de decisiones.
\end{enumerate}

\section*{Requerimientos mínimos del proyecto}
\begin{itemize}
    \item Describir el contexto de negocio y su relevancia.
    \item Plantear al menos cinco (5) objetivos de analítica y visualización.
    \item Utilizar al menos tres (3) conjuntos de datos de origen.
    \item Generar el modelo de datos.
    \item Diseñar y presentar un (1) reporte ejecutivo en PDF que consolide las estadísticas de los indicadores.
    \item Diseñar e implementar al menos una (1) visualización dinámica (dashboard) que incluya:
    \begin{itemize}
        \item Mínimo una pantalla por objetivo.
        \item Mínimo cuatro (4) filtros dinámicos por categorías.
    \end{itemize}
    \item Realizar una visualización análoga en vivo sobre el tema del proyecto e integrarlo como otro insumo de análisis.
    \item La presentación debe demostrar diseño y organización de la información para una comunicación clara y debe contar una historia (storytelling) que guíe a la audiencia desde el planteamiento del problema hasta las recomendaciones.
\end{itemize}

\section*{Entregables}

\subsection*{Primera Entrega – Avance del Proyecto (10\%)}
\begin{itemize}
    \item Descripción del sector y la organización seleccionada.
    \item Definición del contexto y planteamiento del problema.
    \item Preguntas generadoras relacionadas con los objetivos estratégicos.
    \item Objetivos del proyecto y objetivos de analítica/visualización.
    \item Identificación de datasets a utilizar (mínimo tres fuentes).
    \item Modelo lógico de datos (ej. diagrama entidad–relación o dimensional).
\end{itemize}

\subsection*{Segunda Entrega – Proyecto Final (15\%)}
\begin{itemize}
    \item Ajustes y mejoras derivados de la retroalimentación de la primera entrega.
    \item Documento completo del proyecto con:
    \begin{itemize}
        \item Introducción
        \item Contexto
        \item Problema
        \item Objetivos (proyecto y analítica)
        \item Preguntas generadoras.
        \item Descripción de datasets y procesos de preparación/limpieza.
        \item Modelo lógico de datos final.
        \item Registro de la visualización análoga en vivo (fotografías, notas, resultados, reflexión metodológica sobre su utilidad y limitaciones).
        \item Dashboard dinámico con mínimo una pantalla por objetivo y al menos 4 filtros por categorías.
        \item Reporte ejecutivo en PDF con indicadores clave y estadísticas principales.
        \item Conclusiones y Recomendaciones
        \item Lecciones aprendidas.
    \end{itemize}
    \item Archivos anexos: datasets, modelo, y visualizadores.
\end{itemize}

\subsection*{Presentación (10\%)}
\begin{itemize}
    \item Diseño visual y coherencia estética.
    \item Equilibrio entre texto, imágenes y gráficos.
    \item Discurso claro y apropiación del tema por parte del equipo.
    \item Narrativa estructurada: la presentación debe contar una historia (introducción del problema $\rightarrow$ análisis / evidencia $\rightarrow$ conclusiones y recomendaciones).
    \item Organización de la información: jerarquía visual, uso de gráficas y transiciones lógicas que faciliten la comprensión.
\end{itemize}

\section*{Rúbrica de Evaluación}

\subsection*{Primera Entrega – Avance del Proyecto}
\begin{table}[h!]
\centering
\begin{tabular}{|p{4cm}|p{10cm}|}
\hline
\textbf{Criterio} & \textbf{Descripción} \\
\hline
Contexto y sector & Explica de manera clara el sector y la organización seleccionada, mostrando pertinencia y relevancia. \\
\hline
Problema & Plantea un problema bien definido y alineado con el contexto de la organización. \\
\hline
Preguntas generadoras & Formula preguntas útiles que orientan el análisis y la visualización. \\
\hline
Objetivos & Define objetivos claros, alcanzables y coherentes con el problema planteado. \\
\hline
Datasets & Identifica al menos tres fuentes de datos, justificando su relevancia para el análisis. \\
\hline
Modelo lógico de datos & Presenta un diagrama entidad-relación o dimensional coherente con los objetivos y datasets. \\
\hline
Redacción y organización & Entrega un documento estructurado, bien redactado y claro. \\
\hline
\end{tabular}
\end{table}

\subsection*{Segunda Entrega - Proyecto Final}
\begin{table}[h!]
\centering
\begin{tabular}{|p{4.5cm}|p{10cm}|}
\hline
\textbf{Criterio} & \textbf{Descripción} \\
\hline
Ajustes de retroalimentación & Incorpora de manera adecuada los cambios y sugerencias recibidos en la primera entrega. \\
\hline
Contexto y objetivos & Presenta introducción, problema y objetivos (proyecto y analítica) de forma clara y completa. \\
\hline
Preguntas generadoras & Preguntas precisas que guían el análisis y se conectan con los objetivos. \\
\hline
Datasets y preparación & Describe los datasets, procesos de limpieza y preparación de datos de forma rigurosa. \\
\hline
Modelo lógico final & Modelo lógico de datos ajustado, completo y coherente. \\
\hline
Visualización análoga en vivo & Incluye evidencias (fotos, notas, resultados) y reflexión metodológica de la actividad de recolección de datos. \\
\hline
Dashboard dinámico & Presenta un dashboard con mínimo una pantalla por objetivo y al menos cuatro filtros por categorías. \\
\hline
Reporte ejecutivo & Reporte en PDF con indicadores clave, estadísticas relevantes y diseño claro. \\
\hline
Storytelling & Construcción de narrativa que guía desde el problema hasta conclusiones y recomendaciones. \\
\hline
Conclusiones y recomendaciones & Conclusiones sólidas y recomendaciones pertinentes para la organización. \\
\hline
Reflexión crítica & Presenta lecciones aprendidas y reflexión sobre el proceso. \\
\hline
Archivos anexos & Incluye datasets, modelos y visualizadores como parte de la entrega. \\
\hline
\end{tabular}
\end{table}

\subsection*{Presentación Final}
\begin{table}[h!]
\centering
\begin{tabular}{|p{4cm}|p{10cm}|}
\hline
\textbf{Criterio} & \textbf{Descripción} \\
\hline
Diseño visual & Presentación estética, con claridad y equilibrio entre texto, imágenes y gráficas. \\
\hline
Discurso oral & Claridad y apropiación del tema por parte del equipo durante la exposición. \\
\hline
Narrativa & Secuencia lógica: problema $\rightarrow$ análisis/evidencia $\rightarrow$ conclusiones/recomendaciones. \\
\hline
Organización de la información & Jerarquía clara de ideas, transiciones adecuadas y gráficas que facilitan la comprensión. \\
\hline
\end{tabular}
\end{table}

\end{document}